\documentclass[11pt,a4paper]{article}
\usepackage[T1]{fontenc}
\usepackage[utf8]{inputenc}
\usepackage{lmodern,amsmath,amssymb,amsthm,microtype}
\usepackage[margin=28mm]{geometry}
\usepackage[dvipsnames]{xcolor}
\usepackage[colorlinks=true,urlcolor=MidnightBlue,citecolor=MidnightBlue,linkcolor=MidnightBlue]{hyperref}
\usepackage{enumitem}
\setlist{itemsep=3pt,topsep=5pt}
\newtheorem{theorem}{Theorem}
\newtheorem{lemma}[theorem]{Lemma}
\newtheorem{proposition}[theorem]{Proposition}
\theoremstyle{remark}\newtheorem*{remark}{Remark}
\newcommand{\rk}{\operatorname{rk}}
\newcommand{\M}{\mathcal M}
\newcommand{\gen}[1]{\langle #1\rangle}
\title{Minimum-rank overgroups of arbitrary\protect\\ subgroups of free groups\protect\\[5pt]\large A proposed answer to Kourovka Problem 20.11}
\author{AI-generated mathematical note}
\date{Version 1 --- 6 November 2026}
\begin{document}
\maketitle
\begin{center}\small\textbf{Candidate solution --- review pending.}\end{center}
\begin{abstract}
Let $F\le H$, where $H$ is free and $F$ is contained in a finitely generated subgroup of $H$. Let $r$ be the least rank of such an overgroup. We deduce from results of Jaikin-Zapirain and Takahasi that the rank-$r$ overgroups of $F$ form a finite nonempty sublattice of the subgroup lattice of $H$. They therefore have a least and a greatest member, and ordinary intersection and subgroup join preserve minimum rank. This gives affirmative answers to all four assertions in Kourovka Problem 20.11, including when $F$ is infinitely generated. The finitely generated case is already covered by Jaikin-Zapirain's theorem; the purpose here is to give the reduction needed for the unrestricted statement. Priority is not established by this note.
\end{abstract}

\section{The question and the existing theorems}
Problem 20.11 of the Kourovka Notebook, proposed by G.~M.~Bergman, has the following setting. A subgroup $F$ of a free group $H$ is contained in some finite-rank subgroup of $H$. Define
\[
r=\min\{\rk(G): F\le G\le H,\ \rk(G)<\infty\},\qquad
\M(F,H)=\{G:F\le G\le H,\ \rk(G)=r\}.
\]
Here rank is free-basis rank and $\rk(1)=0$. Neither $F$ nor $H$ is assumed finitely generated. The four assertions to decide are:
\begin{enumerate}[label=(\roman*)]
\item $\M(F,H)$ has a greatest member;
\item[$(\mathrm{i}')$] $A,B\in\M(F,H)$ implies $\gen{A,B}\in\M(F,H)$;
\item[$(\mathrm{ii})$] $\M(F,H)$ has a least member;
\item[$(\mathrm{ii}')$] $A,B\in\M(F,H)$ implies $A\cap B\in\M(F,H)$.
\end{enumerate}
A greatest or least member is comparable to every member; mere maximality or minimality is insufficient.

A finitely generated subgroup $A$ of a free group $L$ is \emph{compressed} if each finitely generated subgroup between $A$ and $L$ has rank at least $\rk(A)$. It is \emph{inert} if $\rk(A\cap B)\le\rk(B)$ for every finitely generated $B\le L$.

We use three existing results:
\begin{enumerate}[label=\textbf{\arabic*.}]
\item Compressed finitely generated subgroups of free groups are inert \cite[Corollary 1.2]{JZ}.
\item For a finitely generated initial subgroup of a free group, its minimum-rank overgroups are closed under intersection and subgroup join \cite[Corollary 1.5]{JZ}.
\item A finitely generated subgroup has only finitely many algebraic overgroups in a fixed free group; an overgroup is algebraic when none of its proper free factors contains the initial subgroup. This is Takahasi's theorem, as recorded in \cite[Proposition 2.3]{JZ}.
\end{enumerate}
These are essential external inputs. This note does not offer a new proof of them. In particular, the finitely generated case of the pairwise assertions is already result 2.

\section{Finite intersections}
Write $\M=\M(F,H)$. This family is nonempty by the definition of $r$.
\begin{lemma}\label{lem:finite}
Every nonempty finite intersection of members of $\M$ is again a member of $\M$.
\end{lemma}
\begin{proof}
Take $A,B\in\M$ and set $L=\gen{A,B}$. The group $L$ is free and finitely generated. Any finitely generated overgroup of $A$ inside $L$ contains $F$, and hence has rank at least $r=\rk(A)$. Thus $A$ is compressed in $L$, so result 1 gives
\[
\rk(A\cap B)\le\rk(B)=r.
\]
The intersection is therefore finitely generated. It contains $F$, so minimality of $r$ yields the reverse inequality. Thus $A\cap B\in\M$, and induction proves the lemma.
\end{proof}
This already proves assertion $(\mathrm{ii}')$, without assuming $F$ finitely generated.

\section{The least overgroup}
The main point is to justify the intersection of the entire family. Rank is not generally monotone under subgroup inclusion, so a separate argument is necessary.

\begin{lemma}\label{lem:least}
The subgroup $C=\bigcap_{G\in\M}G$ has rank $r$. In particular, it is finitely generated and is the least member of $\M$.
\end{lemma}
\begin{proof}
Choose $G_0\in\M$. If $r=0$, then $F=G_0=1$ and $\M=\{1\}$, so the conclusion is immediate. Assume $r\ge1$ and choose a free basis $X$ of $G_0$.

For $K\le G_0$, let $\Gamma_K$ be the Schreier graph with vertices the right cosets $Kw$ and, for each $x\in X$, an edge $Kw\to Kwx$ labelled $x$. Inverse traversals are labelled $x^{-1}$. This connected graph is a covering of the rose with basis $X$. Its based fundamental group is $K$, and its first Betti number equals $\rk(K)$.

Let $Q$ be any finite connected subgraph of $\Gamma_C$. Choose representatives $w_1,\ldots,w_m\in G_0$ for its distinct vertices. For $i\ne j$ we have $w_iw_j^{-1}\notin C$. Thus there exists $A_{ij}\in\M$ such that $w_iw_j^{-1}\notin A_{ij}$. Put
\[
B=G_0\cap\bigcap_{i\ne j}A_{ij},
\]
with $B=G_0$ if there are no pairs. Lemma~\ref{lem:finite} gives $B\in\M$ and $\rk(B)=r$. Since $C\le B\le G_0$, there is a canonical labelled graph map
\[
\Gamma_C\longrightarrow\Gamma_B,\qquad Cw\longmapsto Bw.
\]
The exclusions above make this map injective on the vertices of $Q$. An outgoing edge of a specified basis label is unique at each vertex of a Schreier graph. Therefore the map is also injective on the edges of $Q$. Loops and parallel edges cause no exception, since labels and orientations are preserved. Hence $Q$ embeds as a subgraph of $\Gamma_B$.

The cycle rank of a subgraph cannot exceed that of the whole graph. Consequently
\[
b_1(Q)\le b_1(\Gamma_B)=r.
\]
The rank of any connected graph is the supremum of the cycle ranks of its finite connected subgraphs. To see this, choose a spanning tree: each edge outside it contributes an independent cycle, and any finite collection of those cycles lies in a finite connected subgraph. Applying this to $\Gamma_C$ yields $\rk(C)\le r$.

Thus $C$ is finitely generated. Since $F\le C\le H$, minimality of $r$ now gives $\rk(C)\ge r$. Equality follows, and the definition of $C$ makes it the least member of $\M$.
\end{proof}
The argument uses only finite sets of witnesses at a time; no countability assumption on the family $\M$ is needed. The ambient graph is built in the finite-rank group $G_0$, so the rank of $H$ also imposes no restriction.

\section{Joins and the greatest overgroup}
\begin{theorem}\label{thm:main}
For every $F\le H$ in the setting of Problem 20.11, the family $\M(F,H)$ is finite and nonempty, and is closed under ordinary intersection and subgroup join. It has a least and a greatest member. Thus all four assertions in the problem hold.
\end{theorem}
\begin{proof}
Let $C$ be the subgroup from Lemma~\ref{lem:least}. Since $F\le C$ and $C$ lies in every member of $\M$, we have the exact identity
\[
\M=\{G:C\le G\le H,\ \rk(G)=r\}.
\]
The minimum overgroup rank for $C$ is $r$: the group $C$ itself has rank $r$, and any smaller-rank overgroup would also contain $F$.

Take $A,B\in\M$ and again put $L=\gen{A,B}$. Inside $L$, the finitely generated initial subgroup $C$ has minimum overgroup rank $r$. Indeed $A$ is an overgroup of that rank, and a smaller-rank overgroup would contain $F$. Result 2, applied to $C\le L$, now shows that $\gen{A,B}\in\M$. This proves $(\mathrm{i}')$ using only a finite-rank ambient group for the cited theorem.

For each $G\in\M$, the extension $C\le G$ is algebraic. Otherwise a proper free factor $D$ of $G$ would contain $C$. Since $G$ has finite rank $r$, such a free factor has rank strictly less than $r$, contradicting $F\le C\le D$.

Result 3 therefore shows that $\M$ is finite. This application also allows arbitrary ambient free rank. Alternatively choose a finite-rank free factor $P$ of $H$ containing $C$. For an algebraic overgroup $G$, the subgroup $G\cap P$ is a free factor of $G$ by the subgroup theorem for free products. Algebraicity forces $G\le P$, reducing to the finite-rank version of result 3.

Write $\M=\{G_1,\ldots,G_t\}$, where $t\ge1$. Repeated pairwise join closure gives
\[
J=\gen{G_1,\ldots,G_t}\in\M.
\]
It contains every member, so it is the greatest member, proving $(\mathrm{i})$. Lemmas~\ref{lem:finite} and \ref{lem:least} give $(\mathrm{ii}')$ and $(\mathrm{ii})$. Together these establish the stated finite sublattice property.
\end{proof}

\section{Scope, attribution, and verification}
The proposed answer includes infinite-rank $F$ whenever it is contained in a finite-rank overgroup, as well as arbitrary free $H$. The cases $r=0$ and $r=1$ are included. No algorithm for an unspecified infinite subgroup is asserted.

The November 2026 Notebook entry is unmarked and contains no solution comment. That editorial observation and an unsuccessful targeted search do not establish novelty. The central rank-preservation theorem is Jaikin-Zapirain's. The contribution of this exposition is the explicit arbitrary-intersection reduction and its application to the exact four-part formulation. It may be a standard consequence already known to specialists. The arbitrary-intersection principle for inert subgroups is already noted in Kourovka 19.103; the Schreier-graph proof here makes that step explicit.

This manuscript was generated by ChatGPT/Codex. It does not attribute the argument to any human author. Its proof was self-checked against the original quantifiers and the cited theorem statements. A separate Codex reviewer found no mathematical gap in the archived deduction, conditional on the cited results; no repair was needed. The audit and its exact scope are recorded in the accompanying \texttt{review-status-v1.md}; AI review is not human acceptance. The deep proof of the external subgroup-rigidity theorem is not independently re-proved here. No Lean formalization or machine-checked mathematical certificate is supplied.

The candidate is preserved in the research notebook as claim C1 of run \texttt{research-}\allowbreak\texttt{20261106-}\allowbreak\texttt{minrank-}\allowbreak\texttt{rdelaby-7e2c91}, at commit \texttt{faaf280676b2}. The accompanying review record gives the complete identifiers and hashes. Repository validation checks file consistency, not mathematical correctness. No public conversation link was supplied for this version.

\begin{thebibliography}{9}
\bibitem{Kourovka} E.~I.~Khukhro and V.~D.~Mazurov (eds.), \emph{The Kourovka Notebook: Unsolved Problems in Group Theory}, No.~21, November 2026 update, Problem 20.11, p.~158.\par\url{https://kourovkanotebookorg.wordpress.com/wp-content/uploads/2026/11/21tkt.pdf}
\bibitem{JZ} A.~Jaikin-Zapirain, \emph{Free groups are $L^2$-subgroup rigid}, arXiv:2403.09515v2 (2026), Corollaries 1.2 and 1.5, Proposition 2.3, and Section 3.7.\par\url{https://arxiv.org/pdf/2403.09515v2}
\end{thebibliography}
\end{document}
